\input{./._relpath-to-latexroot.ltx}
\documentclass[\latexroot/\projectname]{subfiles}
\input{\latexroot/@resources/subfile-setup.ltx}

\begin{document}

\section{Model}\whenintegrated{\label{model}}
\whenintegrated{\label{sec:model}} % integrated doc: SST for labels

Decompose total government debt $d$ between components owned by banks ($b$) and non-banks ($n$), at home ($h$) and abroad ($f$)
\begin{eqnarray}
    d&=&d_{n}+d_{b}, \\
    d_n &=& d_{h}+d_{f}.
\end{eqnarray}

We assume that the government does not default on home banks, because it is very costly to default on home commercial banks and it does not bring any benefit to default on the home central bank.

The government can default on home non-bank and foreign creditors, $d_h$ and $d_f$. There are different assumptions one can make regarding default.

\medskip

\textbf{Selective default.} First, one can assume that the government can default separately on home and foreign creditors.\footnote{Selective default requires market segmentation at the time of default (see \cite{Broneretal2010}).} Let us assume that the cost of defaulting on debt of type $\theta = h,f$ is the sum of a fixed cost, $\gamma_{\theta}$, plus a variable component proportional to the size of debt that is defaulted on, $\lambda_{\theta} d_{\theta}$. The benefit of defaulting is $\beta_{\theta} d_{\theta}$. The no-default constraint on debt is
\begin{equation}
   \beta_{\theta} d_{\theta} \leq \gamma_{\theta}+\lambda_\theta d_{\theta},
\end{equation}
implying an upper bound on the level of debt,
\begin{equation} \label{eq:nodef1}
    d_{\theta} \leq D_{\theta} \;\;\; \text{with} \; D_{\theta} \equiv \frac{\gamma_{\theta}}{\beta_{\theta}-\lambda_{\theta}}, \;\; \theta=h,f.
\end{equation}
On theoretical grounds, defaulting on residents should be less beneficial than defaulting on foreigners ($\beta_h \leq \beta_f$).\footnote{The sovereign debt literature typically assumes that defaulting on residents has no benefit. However, it could reduce the distortionary taxation required to repay the debt.} Defaulting on residents should be more costly than defaulting on foreigners because it causes more financial disruption at home ($\gamma_h \geq \gamma_f$ and $\lambda_h \geq \lambda_f$). This implies that the government should be able to sell more debt at home, $D_{h} \geq D_{f}$.

Condition (\ref{eq:nodef1}) defines a rectangular region in the $(d_{h}, d_f)$ space (see Figure~\ref{Fig_selective}). For a given total non-homebank government debt $d_n$, the debt composition can be anywhere on the segment AB. The maximum level of $d_n$ is reached when the $d_n = d_h+d_f$ line goes through the angle C, i.e.,  $d_h =D_{h} $ and $d_f =D_{f} $ and the maximum value of $d_n$ is
\begin{equation} \label{eq:Dn}
    D_n = D_{h}+D_{f}.
\end{equation}

\medskip

\textbf{Comprehensive default.} Selective default may be impossible because of market integration or because of \textit{pari passu} clauses in debt contracts. We now show that under some conditions, this does not lead to a very different model of the determination of the composition of debt relative to the case with selective default.

Let us assume that the government either defaults on both $d_h$ and $d_f$ or does not default at all. The benefit of defaulting is $\beta_h d_h + \beta_f d_f$ and the cost of defaulting is $\gamma + \lambda_h d_h + \lambda_f d_f$. Thus, the government does not default if and only if
\begin{equation} \label{eq:nodef2}
    \frac{d_h}{D_h}+\frac{d_f}{D_f} \leq 1,
\end{equation}
where $D_h$ and $D_f$ are defined as before, setting $\gamma_h=\gamma_f=\gamma$.

Observe that the default region is no longer a rectangle but a triangle (see Figure~\ref{Fig_comprehensive}). The maximum level of debt is reached when all the debt is sold at home if $D_h > D_f$. More generally, equation (\ref{eq:Dn}) is replaced by $D_n = \max(D_h, D_f)$.

It would make sense, however, to assume that the no-default region is \textit{strictly} convex. For example, if there are coordination problems between home and foreign creditors, the default cost would be higher when those creditors have more equal weights in the holding of debt. The cost of defaulting could be a concave (rather than linear) function of the debt levels $d_h$ and $d_f$.\footnote{
This intuition can be formalized by allowing default costs to include an interaction term across creditor groups. Let default costs take the form
\( C(d_h,d_f)=\gamma+\lambda_h d_h+\lambda_f d_f+f(d_h,d_f) \),
with \( \partial^2 f / \partial d_h \partial d_f > 0 \).
Such interaction costs imply higher default costs when debt is held more evenly across home and foreign creditors, generating a strictly convex no-default region as in equation~\eqref{eq:nodef3}.
}
In reduced form, convexity could be captured by assuming that (\ref{eq:nodef2}) is replaced by
\begin{equation} \label{eq:nodef3}
    \left( \frac{d_h}{D_h} \right)^\rho+\left( \frac{d_f}{D_f} \right)^\rho\leq 1,
\end{equation}
with $\rho \geq 1$. If there is sufficient convexity, the maximum level of debt $D_n$ is achieved for an interior solution where both residents and foreigners are lending to the government, like in the case of selective default (Figure~\ref{Fig_comprehensive}). As $\rho$ goes to infinity we converge to the case with selective default.

What matters for our results is that there is enough convexity in the no-default region to ensure an interior solution. To simplify the analysis, we will assume that this region is a rectangle, like in the case of selective default, but selective default is not a maintained assumption.

\medskip

\textbf{Financial repression.} The government does not default on banks but accumulating government debt in the banking sector might have costs, for example by crowding out lending to the real sector or because very large bank balance sheets complicate the operation of monetary policy or the management of financial risks.

We assume that the costs of financial repression are zero if bank-held debt $d_b$ is below a threshold $D_b$. Above that threshold, financial repression has a cost $C(d_b)$ where function $C(\cdot)$ is increasing and convex.

\medskip

\textbf{Equilibrium.} Assuming no default, for a given level of total debt $d=d_h+d_f+d_b$, the optimal debt composition minimizes $C(d_b)$ conditional on $d_h \leq D_h$ and $d_f \leq D_f$. If the cost of financial repression is very high, the government might prefer to default but we assume this possibility away as our empirical sample does not include any case of default.

The equilibrium is then easy to determine. Let us define the threshold $D \equiv D_n+D_b$. If $d \leq D$, the government can choose any debt composition $(d_{\theta})_{\theta=h, f, b}$ such that $d_{\theta} \leq D_{\theta}$. If $d > D$, the debt composition is uniquely determined and given by $d_h = D_h$, $d_f = D_f$ and $d_b=d-D_n$.

% Placeholder labels for HAFiscal sections that still reference these (Comparing-policies, Figures)
\whenintegrated{\label{sec:recessions}}
\whenintegrated{\label{sec:ADeffects}}

\begin{figure}[h!]
\centering
\includegraphics[width=0.7\linewidth]{\latexroot/Figures/Fig_dhdfselective}
\caption{No default region: selective default}
\label{Fig_selective}
\end{figure}

\begin{figure}[h!]
\centering
\includegraphics[width=0.7\linewidth]{\latexroot/Figures/Fig_dhdfcomprehensive}
\caption{No default region: comprehensive default}
\label{Fig_comprehensive}
\end{figure}

% Smart bibliography: Only include bibliography if standalone AND has citations
\smartbib

\end{document}
